% Generated draft: requires mathematical and source review.
\documentclass[11pt]{article}
\usepackage[margin=25mm]{geometry}
\usepackage{fontspec}
\setmainfont{Latin Modern Roman}
\usepackage{amsmath,amssymb,mathtools,unicode-math}
\setmathfont{Latin Modern Math}
\usepackage{enumitem,needspace,xurl,hyperref}
\setlength{\parindent}{0pt}
\setlength{\parskip}{5pt}
\setlength{\emergencystretch}{3em}
\newcommand{\N}{\mathbb{N}}
\newcommand{\Z}{\mathbb{Z}}
\newcommand{\Q}{\mathbb{Q}}
\newcommand{\R}{\mathbb{R}}
\newcommand{\C}{\mathbb{C}}
\newcommand{\F}{\mathbb{F}}
\newcommand{\A}{\mathbb{A}}
\newcommand{\PP}{\mathbb{P}}
\newcommand{\E}{\mathbb{E}}
\newcommand{\eps}{\varepsilon}
\newcommand{\dd}{\,\mathrm{d}}
\begin{document}
\section*{UnsolvedMath ID: 30001209}
\textbf{Source problem code:} catalogue entry
\subsection{Definitions and mathematical statement}
A substitution on a finite nonempty alphabet \(A\) is a map \(\sigma:A\to A^{+}\), where \(A^{+}\) consists of nonempty finite words. Its incidence matrix is \(M_\sigma=(m_{ab})_{a,b\in A}\), where \(m_{ab}\) counts occurrences of \(a\) in \(\sigma(b)\). The substitution is primitive if some positive integer power of \(M_\sigma\) has every entry strictly positive. A Pisot number is a real algebraic integer \(\lambda>1\) all of whose other algebraic conjugates have modulus strictly less than \(1\). The supplied record does not specify the convention for a \emph{Pisot substitution}: in particular, it does not say whether having a Pisot Perron eigenvalue suffices or whether an irreducibility or other additional hypothesis is required.

The source asks exactly: ``Is a Pisot substitution determined by its first three correlations? More generally, is every primitive substitution determined by finitely many correlations?''

To display the reconstruction quantifiers without supplying missing conventions, let \(\mathcal C\) be a specified comparison class of substitution objects, let \(\sim\) be the intended equivalence relation on those objects, and let \(D_{\leq k}(\sigma)\) denote their specified correlation data through order \(k\). Once these choices are fixed, the first question can be expressed as
\[
\forall\sigma\in\mathcal P\;\forall\tau\in\mathcal C,\qquad
D_{\leq3}(\sigma)=D_{\leq3}(\tau)
\ \Longrightarrow\ \sigma\sim\tau,
\]
where \(\mathcal P\subseteq\mathcal C\) is the intended class of Pisot substitutions. A pointwise finite-order interpretation of the second question is
\[
\forall\sigma\in\mathcal C_{\mathrm{prim}}\;\exists k\geq1\;\forall\tau\in\mathcal C_{\mathrm{prim}},\qquad
D_{\leq k}(\sigma)=D_{\leq k}(\tau)
\ \Longrightarrow\ \sigma\sim\tau,
\]
where \(\mathcal C_{\mathrm{prim}}\) is the chosen class of primitive substitutions. This permits \(k\) to depend on \(\sigma\); the supplied wording does not specify whether a uniform bound is intended or whether competing objects may lie in a larger class.

For example, under a symbolic, letter-labelled convention, an \(r\)-point correlation for a shift-invariant probability measure \(\mu\) on a two-sided symbolic system \(X\subseteq A^{\Z}\) would be
\[
c^{(r)}_\mu(a_1,\ldots,a_r;n_1,\ldots,n_{r-1})
=\mu\bigl(\{x\in X:x_0=a_1,\ x_{n_j}=a_{j+1}\text{ for }1\leq j<r\}\bigr),
\]
with \(a_i\in A\) and \(n_j\in\Z\). Under this convention, \(D_{\leq k}\) could collect these values for \(1\leq r\leq k\). This is an illustrative convention only: the record does not specify symbolic versus geometric correlations, colour or weight information, displacement restrictions, or the associated system and measure. It also does not specify whether determination concerns the substitution rule, an associated dynamical system, a measure, or a geometric realization, nor the equivalence allowed. Consequently the displayed reconstruction questions are conditional templates, and a definitive mathematical formulation requires these unresolved choices.
\subsection{Short English statement}
The supplied question asks whether three correlation orders determine a Pisot substitution and whether finitely many orders determine each primitive substitution. The meaning of correlation and the object and equivalence to be reconstructed are unspecified, so the formulas retain these choices explicitly. The supplied open status is a dated curation label, not an independently verified conclusion.
\subsection{Sources}
\begin{itemize}
\item Supplied local UnsolvedMath registry record, permanent ID \(30001209\), source problem code \texttt{catalogue entry}, title ``Finite Correlation Determination of Primitive Substitutions.'' The record labels the statement exact and internally readable; this is supplied registry metadata, not independent source verification. No individual proposer is supplied.
\item \emph{Mini-Workshop: The Pisot Conjecture --- From Substitution Dynamical Systems to Rauzy Fractals and Meyer Sets} (2009), \emph{Oberwolfach Reports}, pp.\ 725--766, DOI \texttt{10.4171/owr/2009/13}. Supplied report provenance: \url{https://doi.org/10.4171/owr/2009/13}.
\item Supplied literature source labelled ``Baake--Lenz--Moody report text'': \url{https://ems.press/content/serial-article-files/46214?nt=1}.
\item Supplied literature source labelled ``Lee--Moody--Solomyak primary paper'': \url{https://doi.org/10.1007/s00454-003-0781-z}.
\item Literature assessment dated exactly \texttt{2026-08-21}: the supplied curation labels the question open, says the three-correlation Pisot question is explicitly posed as open, and qualifies the cited correlation theorems as covering pure-point or model-set settings rather than the full Pisot question. This is supplied dated curation derived from the Oberwolfach report, not independently verified here; time-sensitive status and cited claims require independent checking.
\end{itemize}
\end{document}
